\section{Harmonious labellings and harmonious groups}\label{sec:harmonious}

\topic{Additive bases}

Graham and Sloane~\cite{gs80} came to sum labellings from additive number theory.

\begin{problem}[Graham--Sloane~\cite{gs80}, \S1]\label{prob:additive}
Let $n_v(k)$ be the largest $n$ such that some $k$-set $A\subseteq\Z_n$ represents every element of $\Z_n$ as $a_i+a_j$ with $i<j$, and let $v_v(k)$ be the smallest $v$ such that some $k$-set $A\subseteq\Z_v$ has all sums $a_i+a_j$, $i<j$, distinct. Determine $n_v(k)$ and $v_v(k)$.
\end{problem}

For example (from the abstract of~\cite{gs80}), for $A=\{0,1,2,4,8,13,18,22\}$ the $28$ sums $a_i+a_j$ cover $\Z_{24}$, so $n_v(8)\geq24$; in fact $n_v(8)=24$. The motivation was coding theory: ``our interest in $v_v$ stems from its application to error-correcting codes'', since sets with distinct sums modulo $v$ give large constant-weight binary codes $A(k,6,w)$~\cite[\S1]{gs80}.

\topic{Harmonious graphs}

\begin{definition}[Graham--Sloane~\cite{gs80}, \S3]\label{def:harmonious}
A connected graph with $e$ edges is \term{harmonious} if its vertices can be labelled by distinct elements of $\Z_e$ so that the edge sums $f(x)+f(y)$ are distinct. For a tree, where $e=|V|-1$, exactly one label is repeated.
\end{definition}

The edge sums are then exactly $0,1,\dots,e-1$. A harmonious labelling of $K_v$ is exactly a set attaining $n_v(v)=v_v(v)=\binom v2$, and hence $K_v$ is harmonious if and only if $v\leq4$~\cite[Theorem~7]{gs80}; for instance $\{0,1,2,4\}$ works for $K_4$. Almost all graphs are not harmonious, while odd cycles, ladders and wheels are~\cite{gs80}. Figure~\ref{fig:harmonious-path} shows a path with $3$ edges labelled in $\Z_3$.

\begin{figure}[ht]
\centering
\begin{tikzpicture}[every node/.style={font=\small}]
    \foreach \i/\l in {0/0,1/1,2/1,3/2} {
      \node[bvertex] (h\i) at (1.5*\i,0) {}; \node[below=2pt] at (h\i) {$\l$};}
    \draw (h0)--node[above, RainbowRed] {$1$} (h1);
    \draw (h1)--node[above, RainbowRed] {$2$} (h2);
    \draw (h2)--node[above, RainbowRed] {$0$} (h3);
  \end{tikzpicture}

\caption{A path with $3$ edges labelled in $\Z_3$, with the label $1$ repeated; the edge sums are $1,2,0$.}
\label{fig:harmonious-path}
\end{figure}

\begin{theorem}[Graham--Sloane~\cite{gs80}, Theorems~12--13]\label{thm:gs-trees}
Every caterpillar is harmonious, and every tree with at most $9$ vertices is harmonious.
\end{theorem}

\begin{conjecture}[Graham--Sloane~\cite{gs80}, \S5]\label{conj:gs}
Every tree is harmonious.
\end{conjecture}

This has been verified for all trees with at most $31$ vertices (Fang, 2011)~\cite[\S2.1]{gal25}. Graham and Sloane write ``cf.\ [Bloom]'': the conjecture is the harmonious twin of the graceful one. For example, the tree of Figure~\ref{fig:fork}, labelled $0,1,2,1,3$ in $\Z_4$, is harmonious.

\topic{Lengths and directions}

Label the vertices of a regular $m$-gon by $0,\dots,m-1$ (Figure~\ref{fig:chords}). The chord $\{x,y\}$ has the same length as $\{x',y'\}$ if and only if $x-y\equiv\pm(x'-y')$, and it is parallel to $\{x',y'\}$ if and only if $x+y\equiv x'+y'$~\cite[\S3]{gs80}. So a $\rho$-labelling asks that no two edges have equal length, and a harmonious labelling asks that no two edges be parallel.

\begin{figure}[ht]
\centering
\begin{tikzpicture}[every node/.style={font=\footnotesize}]
    \polygonPoints{a}{7}{1.25}{0}
    \foreach \i in {0,...,6} \foreach \j in {0,...,6} {
      \ifnum\i<\j
        \pgfmathtruncatemacro{\d}{min(\j-\i,7-\j+\i)}
        \ifcase\d\or \draw[col1] (a\i)--(a\j); \or \draw[col2] (a\i)--(a\j);
        \or \draw[col3] (a\i)--(a\j); \fi
      \fi}
    \polygonLabels{a}{7}{1.25}{0}
    \node[align=center] at (0,-1.95) {colour $=$ length:\\
      \textcolor{RainbowRed}{$1$}, \textcolor{RainbowBlue}{$2$}, \textcolor{RainbowGreen}{$3$}};
    \polygonPoints{b}{7}{1.25}{5}
    \foreach \i in {0,...,6} \foreach \j in {0,...,6} {
      \ifnum\i<\j \draw[gray!25] (b\i)--(b\j); \fi}
    \draw[col1, line width=1.2pt] (b1)--(b6) (b2)--(b5) (b3)--(b4);
    \draw[col2, line width=1.2pt] (b0)--(b1) (b2)--(b6) (b3)--(b5);
    \polygonLabels{b}{7}{1.25}{5}
    \node[align=center] at (5,-1.95) {colour $=$ direction:\\
      \textcolor{RainbowRed}{$x+y\equiv0$}, \textcolor{RainbowBlue}{$x+y\equiv1$}};
  \end{tikzpicture}

\caption{The chords of a regular $7$-gon, coloured by length (difference up to sign) and by direction (sum modulo $7$).}
\label{fig:chords}
\end{figure}

\topic{Harmonious groups}

\begin{definition}[Beals--Gallian--Headley--Jungreis~\cite{bghj91}]\label{def:harmonious-group}
A finite group $G$ is \term{harmonious} if its elements can be listed $g_1,\dots,g_n$ so that $g_1g_2,\ g_2g_3,\ \dots,\ g_{n-1}g_n,\ g_ng_1$ is again a list of all elements of $G$.
\end{definition}

``Our definition was motivated by the concept of a harmonious graph invented by Graham and Sloane''~\cite[\S1]{bghj91}.

\begin{fact}[Beals et al.~\cite{bghj91}, \S2]\label{fact:harmonious-cm}
$G$ is harmonious if and only if it has a complete mapping that is a single $|G|$-cycle, $g_1\mapsto g_2\mapsto\dots\mapsto g_n\mapsto g_1$.
\end{fact}

For example, $\Z_7$ is harmonious via $0,1,\dots,6$, with sums $1,3,5,0,2,4,6$; $\Z_2\times\Z_4$ is harmonious; $\Z_4$ and $\Z_2^2$ are not. There are two algebraic obstructions.

\begin{theorem}[Paige; Beals et al.~\cite{bghj91}, Theorem~3.1 and Corollary~3.2]\label{thm:harm-paige}
If $g_1,\dots,g_n$ is a harmonious listing, then $g_1g_2\cdots g_n$ lies in the commutator subgroup. Hence a group of even order with a cyclic Sylow $2$-subgroup is not harmonious.
\end{theorem}

This is the Hall--Paige obstruction.

\begin{theorem}[Beals et al.~\cite{bghj91}, Theorem~3.3]\label{thm:harm-z2k}
The elementary abelian $2$-groups $\Z_2^k$ are not harmonious.
\end{theorem}

Indeed, in $\Z_2^k$ the equation $gh=e$ forces $g=h$, so the identity is never the product of two consecutive, distinct terms. For abelian groups these are the only obstructions.

\begin{theorem}[Beals et al.~\cite{bghj91}, Theorem~6.6]\label{thm:harm-abelian}
A finite abelian group is harmonious if and only if its Sylow $2$-subgroup is trivial or non-cyclic and it is not an elementary abelian $2$-group.
\end{theorem}
